% !TeX root = ../../main.tex
\section{Discussion and Limitations}
\label{sec:mobiwac:discussion}

A single model serves both tasks in one forward pass, on next region it outperforms the dedicated
model at Texas and California, the two datasets with the largest region vocabularies, and at the
other four it stays within the two-point margin registered before any result was read. On next
category it outperforms the dedicated model at Florida, and the five remaining differences are not
resolved, each of them bounded within half a point of zero. Throughout, the region output keeps its
own private spatial path (Table~\ref{tab:mobiwac:results}; Figure~\ref{fig:mobiwac:deltas}). The
practical reading is that one model replaces two while staying inside the margin the analysis plan
set as the price worth paying, and does better than two where the region label space is largest.
The four region cells inside that margin are all small deficits rather than ties, the largest of
them $0.87$ Acc@10 points at Alabama, so the trade is a measured one and not a free substitution.

The design gives the two tasks a shared representation while keeping a private spatial path for the
region output. The evidence here does not separate their contributions. Where the trunk was isolated
directly, the arms moved both tasks by small amounts that a screen of that size cannot distinguish
from noise, and those arms were not run at the two datasets that carry the region advantage. The
reading this supports is a cautious one: the shared representation appears able to help, and how much
it helps is a question for further work rather than one this study settles. What can be stated
without that work is narrower, and it is a claim about the model rather than about transfer between
its tasks: this design, shared representation and private path together, produces a joint region
output above two dedicated models at the two datasets with the largest region vocabularies.

The category picture at Texas and California invites a specific hypothesis. Batch size was searched
for the dedicated category model at every dataset and the learning rate at four of them, with both
searched for the joint model at the smaller datasets; at Texas and California the joint model runs a configuration transferred from
those searches rather than one validated on them. Texas carries the second-largest category deficit
of the six ($-0.13$; Alabama's $-0.19$ is the largest, and it is a dataset where the joint model was
searched), and California's is under a hundredth of a point. A search conducted for the joint model
at those two datasets would test whether the transferred configuration is what holds their category
output slightly below the dedicated model. Until it is run this remains a hypothesis rather than an
explanation, and Alabama is a reminder that a searched dataset can show a deficit of the same size.

A service could treat the model's top-ranked next regions as an anticipatory
set: at California, ten regions out of $8{,}501$ contain the true next region $64.54$
percent of the time, and at Texas, ten out of $6{,}553$ contain it $66.15$ percent of the time.
Both figures are the joint model's next-region Acc@10 from
Table~\ref{tab:mobiwac:results}, read here as a shortlist hit rate: a system that surfaced ten
census tracts would be right about two times in three, over a candidate set three orders of
magnitude larger. Where the shortlist misses, the geographic size of the error is the quantity that
would matter to such a service, and measuring it requires the per-visit predictions that the
evaluation path does not retain, so it is left to future work. This remains motivation, not a
measured service result.


Four limits qualify these results. First, the representation is trained once over all places; a
per-fold rebuild from training users only changed the results by at most $0.33$ Acc@10 and $0.29$
macro-F1 across three datasets at one seed, with the qualification on the region half recorded in
Section~\ref{sec:mobiwac:setup-windows}. Second, epoch selection consults the fold that the score is
then read on (Section~\ref{sec:mobiwac:setup-windows}), so every absolute score reported here is optimistic. The
comparison between the joint model and the dedicated models is affected far less, for two reasons that we can state
rather than assume: the selection rule is the same for both models on the same folds, an epoch chosen on validation and read
on that fold, with each model selected on its own validation objective (Section~\ref{sec:mobiwac:results-part2}),
and, for the category comparison, the dedicated model receives the wider search: batch size at all
six datasets and the learning rate at four, against a joint-model search that covers four of the six
and leaves Texas and California with a transferred configuration, while both sides of the region
comparison run a fixed configuration (Section~\ref{sec:mobiwac:setup-windows} states the coverage
per knob and dataset). Where the dedicated search is the wider of the two, the residual favors the
dedicated model, which makes the reported category difference conservative there. At Florida and
California, where the dedicated learning rate was not varied, the two searches are closer in extent
and the reading does not apply; Florida is also the one dataset where the joint model outperforms on
this task. It does not follow that the bias cancels exactly. Third, we do not
build or evaluate a mobility-aware service; it is background motivation, and this chapter's claims are the
prediction results themselves.
Fourth, each
visit node in the representation graph draws on the visits that precede it, so a vector encodes the visit
together with its own history; the graph does not pass information from a later visit back to an earlier
one, in training or at readout, which is what keeps a node from carrying a feature of the target it is
used to predict.

Our representation is a fixed per-visit vector that any downstream model
can consume. Moura et al.~\cite{moura2025mobilityaware}, whose structural
analysis of tourist check-ins reads past visits only, name machine learning
as a next step; this study moves in that direction.
